%% ==========================================================================
%%  SECTION 2 --- WHAT "PRIVATE" MEANS                          [Track T4]
%%
%%  Job: establish that "private recommendation" names several INCOMPARABLE
%%  guarantees, so the numbers in later sections cannot be compared naively.
%%  This section is what stops the survey from being a list. ~1.5 pages.
%% ==========================================================================
\section{What ``Private'' Means}
\label{sec:threat}

\todo{The point of this section: four communities all say ``privacy-preserving
      recommendation'' and mean four different things. A reader who does not
      have this straight will mis-read every comparison in \S4 and \S5.}

\subsection{The four families}
\label{sec:threat:families}

\todo{One short paragraph each, with a representative citation and an honest
      statement of what each does NOT protect.
      (a) MPC / secret sharing --- \pirsonaplain{}, \nudgeplain{}. Strong,
          but assumes non-collusion among a fixed set of servers.
      (b) Homomorphic encryption --- single-server, no non-collusion
          assumption, orders of magnitude slower.
      (c) Federated learning + DP --- no cryptographic guarantee on the
          intermediate model; \nudgeplain{} \S3.1 Remark 3.1 is explicit that
          federated systems leak through logs and intermediate model versions.
      (d) Trusted hardware --- deployed and fast, but trades cryptography for
          hardware trust and is vulnerable to side channels.}

\begin{table}[t]
\centering
\todo{Table: family $\times$ \{trust assumption, what is hidden, what leaks,
      typical cost\}. This table is the section's deliverable --- if a reader
      takes one thing from \S2, it should be this. Use booktabs.}
\caption{Privacy notions for recommendation and what each actually buys.}
\label{tab:threat:families}
\end{table}

\subsection{Corruption models}
\label{sec:threat:corruption}

\todo{Semi-honest versus malicious; honest majority versus dishonest
      majority; static versus adaptive corruption. Define the real/ideal
      paradigm briefly and cite Lindell's ``How to Simulate It''.
      Be precise about the two base papers:
      \nudgeplain{} is semi-honest, honest-majority, 3 servers, tolerates
      compromise of the entire secret state of ONE;
      \pirsonaplain{} is pairwise non-colluding across $s+1$ servers, with a
      PIR layer that is computationally 1-private against a MALICIOUS server.
      That difference matters and is easy to get wrong.}

\subsection{Leakage that survives the protocol}
\label{sec:threat:leakage}

\gap{This subsection plants the seed for \S8. Do not resolve it here.}

\todo{The idea that a protocol can be perfectly secure and still leak, because
      the leakage is in the FUNCTIONALITY rather than the protocol. Two
      independent arrivals at the same point:
      Asharov et al.\ (PETS'18) on genomic search --- an adaptive querier
      learns a great deal even against an IDEAL functionality;
      PICS (eprint 2025/1071) on input inconsistency across executions.
      Then note that \nudgeplain{} publishes $\itemembed$ in the clear by
      design, and ask what that gives an adversary. Leave it hanging.}
